\documentclass[12pt]{article}
\usepackage[margin=25mm]{geometry}
\usepackage{url}
\begin{document}
\pagestyle{empty}

\noindent\today

\vspace{1em}
\noindent Dear Editor,

\vspace{1em}
\noindent We are pleased to submit our manuscript entitled ``Characterizing Sensor Drift and Cross-Calibration in Korea's National Air Quality Monitoring Network'' for consideration as a Technical Design Note in \textit{Measurement Science and Technology}.

This study presents the first network-scale empirical characterization of sensor performance across 350 stations in Korea's AirKorea monitoring network, analyzing 748,792 hourly measurements of six pollutants. Our key contributions include:

\begin{enumerate}
\item Quantification of sensor drift prevalence (33.8--61.3\% of stations) using both linear regression and CUSUM change-point detection, with CUSUM identifying drift onset at a median of 18 days.

\item Spatial cross-calibration analysis revealing that 23.7\% of co-located station pairs exhibit systematic biases exceeding 5~$\mu$g/m$^3$ for PM$_{2.5}$.

\item A novel Sensor Health Index (SHI) that integrates drift, noise, and data availability into a single composite metric for prioritizing maintenance resources.
\end{enumerate}

These findings advance measurement science by providing actionable benchmarks for calibration scheduling, uncertainty quantification, and quality control in national air quality monitoring networks. The proposed SHI offers a practical tool for network operators to identify stations requiring urgent attention.

The manuscript has not been published previously and is not under consideration elsewhere. All data were collected from publicly available sources (AirKorea Open API), and the analysis code is publicly accessible.

Thank you for your consideration.

\vspace{2em}
\noindent Sincerely,\\
Yongjun Kim\\
Department of Software, Ajou University\\
Suwon 16499, South Korea\\
yjkim0123@ajou.ac.kr

\end{document}
